\section{A lower bound of $\Omega(n^{1/3})$ for all constant-length cycles }\label{sec:Ck-lb}

For the ease of readability, we first restate Theorem~\ref{thm:lb_ck_const_alpha}

\medskip
\noindent{\textbf{Theorem~\ref{thm:lb_ck_const_alpha}}}~
\emph{\GenLB}


\medskip

% \commentt{maybe talk about dependence on $k$ somewhere?}

\iffalse
\begin{theorem}\label{thm:lb_ck_const_alpha}
Let %$\eps>0$,   
$k\geq 6$ and  $n\in \mathbb{N}$.
Any one-sided error tester for the property of $C_k$-freeness in graphs with $n$ vertices and arboricity $\alpha=2$, must perform $\Omega(n^{1/3})$ queries for constant $\eps$.
\end{theorem}
\fi

As noted in the introduction, the lower bound stated in Theorem~\ref{thm:lb_ck_const_alpha} also applies to graphs with non-constant arboricity by adding a $C_k$-free subgraph with  higher arboricity. For an odd $k$, it suffices to add a dense bipartite graph, and for even $k$, by the Erd\H{o}s girth conjecture~\cite{E63}, one can add a subgraph with arboricity $n^{2/k}$.

\ForFuture{I still feel a bit weird about this, so just leaving a comment}


To prove the theorem we shall reduce (in a non black-box manner) from the following result:

\begin{theorem}[\protect{\cite[Lemma 2]{AKKR08}}]\label{thm:triangle_freeness}
    Every one-sided error algorithm for testing triangle-freeness in graphs with $\trn$ vertices and average degree $d$ must
perform $\Omega(d,\trn/d)$ queries. This lower bound holds even when %$d_{max}=O(d)$.
the maximum degree is $O(d)$.
\end{theorem}

\begin{lemma}[Implicit in the proof of Lemma 2 in \cite{AKKR08}]\label{lem:tri_free_hard_family}
There exists a family of (multi-)graphs $\mG_{\trn}$ such that 
the following holds.
\begin{itemize}
    \item The graphs  are tripartite and have $\trn$ vertices in each part
    \item The graphs are $d$-regular for $d=\Theta(\sqrt{\trn})$, 
    \item All but at most a small constant fraction of  graphs in the family  are $\Omega(1)$-far from triangle-free. 
    %\commentt{DR: bit a bit more precise, by either stating a constant or saying that all but at most a $O(1)$-fraction...}\talya{is it OK now?}\commentt{added ``small'' and I think that Ok}
\end{itemize}

Every algorithm that uses neighbor
and/or vertex-pair queries on a graph selected uniformly at random from  $\mG_{\trn}$ must perform $\Omega(\sqrt{\trn})$ queries before it 
views a triangle with sufficiently high constant probability.  
Furthermore, this holds even when the algorithm knows in advance which vertex belongs to which of the three parts.
\end{lemma}

% We are now ready to prove Theorem~\ref{thm:lb_ck_const_alpha}.

\begin{proofof}{Theorem~\ref{thm:lb_ck_const_alpha}}
To prove the theorem, assume towards a contradiction that there exists a one-sided error algorithm $\mA$ for testing $C_k$-freeness using at most %$o(N^{1/3})$ 
$n^{1/3}/c$ queries  (for some constant $c$)
in graphs with $n$ vertices and arboricity $2$.
We shall prove that  there exists an algorithm $\trA$ that % tests triangle-freeness 
finds, with high constant probability, a triangle in graphs selected uniformly from $\mG_{\trn}$ using at most $\sqrt{\trn}/c'$ queries 
 for some constant $c'$, thus reaching a contradiction to Lemma~\ref{lem:tri_free_hard_family}.
% \commentt{I think we need to be a bit more precise here regarding why it is a contradiction. Supposedly, a one-sided-error algorithm doesn't necessarily have to find a $C_k$ }

Algorithm $\trA$ will operate as follows. Given query access to a graph $\trG\in \mG_{\trn}$, algorithm $\trA$ (implicitly) defines a graph $G$ and invokes algorithm $\mA$ on it, where each query to $G$ is answered using a constant number of queries to $\trG$ (in expectation). 
Once $\mA$ terminates its execution on $G$, 
% if it rejects, then it must have 
if it found  a $C_k$ in $G$, 
then $\trA$ returns a corresponding triangle in $\trG$, 
% $\trA$ returns the same answer on $G$, 
where there is a small  probability ($o(1)$) that $\trA$ terminates early (before $\mA$ makes its decision). 

We shall show that 
% if $G$ is triangle-free, then $G'$ is $C_k$-free, and 
if $\trG$ is $\treps$-far from triangle-free, then $G$ is $\eps=\Omega(\treps)$-far from $C_k$-free. In such a case (when executed with a distance parameter $\eps$),  $\mA$ must reject $G$ with high constant probability. Since $\mA$ has one-sided error, when it rejects, it must have evidence in the form of a $C_k$.\footnote{This statement holds for the graphs constructed, and assuming the number of queries performed by $\mA$ is indeed at most $n^{1/3}/c$. The reason is that if after performing these many queries the algorithm did not detect a $C_k$, then the algorithm's  ``knowledge graph'' (which includes all queried edges as well as degrees) can be completed to a graph that is $C_k$-free. \ForFuture{work on phrasing}}
% \commentt{Should double check this} \talya{looks like need to use the fact that lower bound is smaller than $n' = n^{2/3}$}} 
% To be precise, this is the case when the number of queries it performs is sublinear in the size of the graph, since if the graph induced by the queries it performed does not contain a $C_k$, then it is consistent with some graph (having constant arboricity) that contains no $C_k$s.
%\commentt{DR: Do we need to justify this more}
% Hence, 
% if $\mA'$ accepts every $G'$ that is $C_k$ free, then $\mA$ accepts every graph $G$
% if $\mA$ is correct on $G$ with some probability $1-\delta$, then $\mA'$ is correct on $G'$ with probability $1-\delta'$, where $\delta' = \delta + o(1)$.
% \commentt{DR: shouldn't it be the opposite? Also, I am now wondering about two things. First (related to a previous comment), we need to say something about the need to see a $C_k$ in order to reject, and the second is whether it suffices to only talk about maintaining distance given that graphs in the family we work with are (almost) all far from triangle-free}


For every graph $\trG$ in $\mG_{\trn}$,   we consider an orientation of its edges so that for  edge $(u,v)$  if $id(u)< id(v)$, 
% \commentt{DR: reminder to add notion of ids to preliminaries}  
then the edge is oriented from $u$ to $v$.  For ease of presentation, assume for now that $k$ is divisible by $3$.
The graph $G$ is constructed in two steps where the first is deterministic and the second is randomized. First a graph $G^+$ is obtained by replacing each oriented edge $e=(u,u')$ in $\trG$ with a path of length $k/3$, where the original endpoints are kept in $G^+$, and $\ell=k/3-1$ new vertices $v_{e}^1, \ldots, v_{e}^{\ell}$ are added along the path (where $u=v_{e}^0$ and $u'=v_e^{\ell+1}$).
We refer to the set of vertices in $G^+$ that originated from endpoints of edges in $\trG$ as ``original vertices'', and to the set of new vertices as ``path vertices''.
See Figure~\ref{fig:gen-lb} for an illustration.
(Note that even if the graph $\trG$ has parallel edge, since every edge in $\trG$ is replaced by a distinct path in $G^+$, the resulting graph $G^+$ has no parallel edges.)

In the second step, the graph $G$ is obtained by applying a random permutation on the ids of the vertices of $G^+$, where the permutation is applied separately to the ids of the original vertices and to the ids of the path vertices. That is, the algorithm $\trA$ knows at the beginning of its execution, what is the set of ids of the original vertices and  of path vertices. If $k$ is not divisible by $3$, then  to obtain $G^+$ from $\trG$, 
a similar process to the above is done with either $\lceil k/3\rceil, \lceil k/3\rceil, \lfloor k/3\rfloor$ or $\lceil k/3\rceil, \lfloor k/3\rfloor, \lfloor k/3\rfloor$, where all edges between the same two parts in the original graph $\trG$, are replaced by a path of the same length in $G^+$.


Recall that all graphs $\trG$ in $\mG_{\trn}$ have $\trm=\Theta(\trn\cdot d)=\Theta(\trn^{3/2})$ edges so that $G$ has $n=\Theta (k\cdot \trm)=\Theta(\trn^{3/2})$ vertices and  $m=\Theta(\trn^{3/2})$ edges (recall that $k$ is a constant). Also note that $G$ has arboricity $2$, as for any subgraph  of $G$, its average degree is at most $2$.  

Observe that by the above construction, every triangle in $\trG$ is transformed into a $k$-cycle in $G$, and every $k$-cycle in $G$ originates from a triangle in $\trG$. Therefore, if a $k$-cycle is detected in $G$, then a triangle is detected in $\trG$.
We also claim that the distance is preserved (up to a constant factor), i.e., that $dist(G, C_k\mbox{-}free)= \Omega(dist(\trG, C_3\mbox{-}free))$. 
% \commentt{DR: reminder to add notation of $dist$ (or $\textrm{dist}$ to preliminaries.}  
To verify this, let $R$ be a 
minimum-size set of edges so that if we remove $R$ from $G$, then it becomes $C_k$-free. Map each edge $e\in R$ that resides on some path in $G$ to the edge $\tre\in E(\trG)$ that this path replaced, and let $\trR$ be the set of edges that $R$ maps to. Then $|\trR|= |R|$, and  removing $\trR$ from $\trG$ will make $\trG$ triangle-free.\footnote{Assume the contrary, i.e., that after removing $\trR$ from $\trG$, there exists a triangle in $\trG$. Then the three paths in $G$ that replaced this triangle are disjoint from $R$, as otherwise this path would not have survived. Therefore, $G$ after removing $R$ is not $C_k$-free which is a contradiction to the definition of $R$.} Therefore,

\begin{equation} dist(\trG, C_3\mbox{-}free)\leq \frac{|\trR|}{\trm}= \frac{|R|}{\trm}\leq \frac{|R|}{m\cdot \lfloor 3/k\rfloor}=\Omega(dist(G,C_k\mbox{-}free))\;,
\end{equation}
%(since $m'=\Theta(m)$). 
as claimed.
% Hence, we get that $\eps'=\Omega(\eps)$ as claimed. \commentt{DR: remind that $k$ is a constant? Do we want $\eps$ to equal the distance, or is the distance greater than $\eps$? In the next paragraph we have $\eps$-far}


Hence, given access to $\trG$ in $\mG_{\trn}$,  we shall invoke $\mA$ on $G$.
If $\trG$ is $\treps$-far from being triangle-free (for some constant $\treps$, as almost all graphs in $\mG_{\trn}$ are), then $G$ is $\treps/c'$-far from being $C_k$-free for some constant $c'$. Therefore, by the assumption on $\mA$, we can, with high constant probability, using %$o((n')^{1/3})=o(\sqrt n)$ 
at most $n^{1/3}/c \leq \sqrt{\trn}/c'$ %\trn^{1/2}/c'$
queries on $G$, 
% answer correctly whether $G$ is triangle-free. 
obtain a triangle in $\trG$.
Thus we get a contradiction to Lemma~\ref{lem:tri_free_hard_family}. 

It remains to show that  the invocation of $\mA$ on $G$ can be simulated by giving $\trA$ query access to $\trG$, with the same order of query complexity. % and success probability.


To answer queries to $G$ using queries to $\trG$, $\trA$ constructs uniform permutations $\sigma_1\in\Pi_{\trn}$ and $\sigma_2 \in \Pi_{n-\trn}$ ``on the fly", so that  $\sigma_1$ maps original vertices in $G$ to original vertices in $G^+$ and $\sigma_2$ maps path vertices in $G$ to path vertices in $G^+$.  
%\commentt{A reminder to say where we use the assumption stated in the lemma that the algorithm may known the partition of labels into the three subsets}

Let $I_{1,t}\subset[\trn]$ and $I_{2,t}\subset [\trn+1, n]$ be the set of indices,  for which $\sigma_1$ and $\sigma_2$, respectively,  were not yet decided on by the $t$th query.  Further let  $\sigma(I_{1,t})$ and $\sigma(I_{2,t})$ be the set of indices in $[\trn]$ and $[\trn+1, n]$, respectively, that were not yet mapped to by $\sigma$ by the $t$th query.
For any $t\geq 0$,  the $t+1$ query is   answered by $\mA$ as follows. 



\paragraph{Degree queries $deg(u)$:} 



Recall that every $\trG\in \mG_{\trn}$ is $d$-regular, so that all the original vertices in $G$ have degree $d$. Also, by construction, all path vertices in $G$ have degree exactly $2$. Finally, by construction, the set of ids of each group of vertices is known in advance. Hence, we can assume without loss of generality that no degree queries are performed.

%We further assume that when $\mA$ answers a neighbor query, answers a pair query affirmatively, or answers a degree query on a path vertex, it returns the  length-$\ell$ path that the revealed edge or path-vertex belongs to. We refer to such answers as a ``path-answer", and we shall give more details below.
%\commentt{DR: I think something in the combination of the previous version and the new one isn't exactly consistent. First, we just said that it doesn't ask degree queries. Second, later on, in pair queries, it is assumed that such queries are not performed on path vertices.}

%\commentt{CONTINUE CHANGE IN NOTATIONS FROM HERE ON}

\paragraph{Neighbor queries $nbr(u,p)$:} 
The query is answered differently depending on whether $u\in [\trn]$ or $u\in [\trn+1,n]$.
If $u\in [\trn]$, then $\trA$ first selects $\sigma_1(u)$ by choosing an index in $[\trn]\setminus \sigma_1(I_t)$ u.a.r.. Then $\trA$ performs a neighbor query on the original vertex $\sigma(u)$ in $\trG$. That is, $\trA$ performs the query $nbr(\sigma(u), p)$ on $\trG$. 
    Let $\tre=(\sigma(u),v)$ denote the returned edge. 
    If $v$ was not previously discovered (i.e., was not previously queried or was returned as an answer to a previous query), then  
    $\trA$ decides on its id in $G$ , i.e., $\sigma^{-1}_1(v)$, by choosing uniformly from the indices in $[\trn]\setminus I_{1,t}$. 
    % \commentt{Early termination otherwise?} 
    $\trA$ also performs a \emph{path labeling  process} on the path that replaces the edge $\tre$ in $G$:
    Let $v_{\tre}^1, \ldots, v_{\tre}^{\ell}$ be the path vertices in $G^+$ that 
    %were added to $G^+$ in 
    on the path that replaced the edge $\tre$ in $\trG$ (where $\sigma(u)=v_{\tre}^0$ and $v=v_{\tre}^{\ell+1}$).
    $\trA$  labels the path vertices  by choosing for each $v_{\tre}^r$ a value $\sigma_2^{-1}(v_{\tre}^r)$ u.a.r. in $[\trn+1, n]\setminus I_{2,t}$.
    Finally, $\trA$ answers the query with all the vertices of the path and their ids. % (and degrees). This is what 
    We refer to this as a ``path answer".

    We turn to the case that the neighbor query is performed on a path vertex. That is, $u\in [\trn+1, n]$.  
    In this case we would like to map $u\in G$ to an edge $\tre\in \trG$.  
    Recall that the id $u$ not only specifies that $u$ is a path vertex, but also to which two parts its path belongs to. Assume without loss of generality that $u$ belongs to the set of path vertices that reside on paths between parts $1$ and $2$. In order to map $u$ to an edge between parts $P_1,P_2$, we need to first sample an edge $e$ between $P_1$ and $P_2$ uniformly at random  (we shortly explain how).  If the edge $\tre\in \trG$ has already been  discovered, then $\trA$ terminates and returns that $\trG$ is triangle free. This only happens with negligible probability $O(Q/m)=o(1)$, where $Q$ is the number of queries performed by $\mA$. Otherwise, $\trA$  draws   a uniform index $r\in [\ell]$,  
  and lets $\sigma(u)=v_{\tre}^{r}$. Then 
    $\trA$ performs a path labeling process on the remaining path vertices and returns a path answer for $\tre$.
    
    To sample a uniform edge in $\trG$ between two specific parts, say $P_1$ and $P_2$,  $\trA$ samples $w$ in the set of indices of $P_1$ and $r'\in [d]$ u.a.r. and performs  neighbor queries $nbr(w,r')$ on $\trG$ until the returned neighbor belongs to the desired part, in this case $P_2$. 
    Since $\trG$ is $d$-regular, this results in a uniformly distributed edge $\tre$ in $E(P_1,P_2)$. Note that since every vertex $v\in G$ has exactly half of its neighbors in each of the other two parts, the above process returns an edge in the desired part with probability $1/2$. Hence, the expected number of queries per one query simulation  is constant.

    
    \paragraph{Pair queries $pair(u_1,u_2)$:} 
    If the two queried vertices are original vertices in $\trG$ then the answer to the query is no. If at least one of them is a path vertex, then $\mA$ maps it to an edge, as described above for neighbor queries, and answers accordingly.
    % already knows the answer to the query since path vertices are always returned along with the $\ell$-path they belong to. Hence,  here too we can assume without loss of generality that no pair queries are performed.

% \paragraph{Final answer.} 
%Once $\mA'$ completes its execution, if it found a $C_k$ in $G'$, then this $C_k$ corresponds to a triangle in $G$ (over the original vertices of $G$ in the $C_k$), and $\mA$ returns this triangle. 
%\commentt{DR: I was going to continue, but wouldn't it be repeating what we already said?}
\iffalse
Once $\mA'$ finishes its execution and returns some answer, $\mA$ returns the same answer. 
\commentt{DR: We need to fix this too. If we are happy with what is written in the introductory text, then we can modify this part accordingly}
  As noted previously, if $G$ is triangle-free, then $G'$ is $C_k$-free. Since $\mA'$ has one-sided error, and upon early termination of $\mA'$, $\mA$ accepts, $\mA$ also has one-sided error. On the other hand, if $G$ is $\eps$-far from triangle-free, then (also as noted previously), $G'$ is $\Omega(\eps)$-far from $C_k$-free. In this case, $\mA$ answers incorrectly either if $\mA'$ answers incorrectly or in the case of early termination. Since the first event happens with probability $\delta$, and the second with probability $o(1)$, we get that $\mA$ answers incorrectly with probability at most $\delta'=\delta=o(1)$. 
\fi

\bigskip
Finally, as each query to $G$ can be simulated in a constant number of queries to $\trG$ in expectation, it follows that with high probability $1-o(1)$ (over the randomness of $\trA$), the total number of queries does not exceed that of $\mA$ by more than a constant factor, %$C'$ for some $C'$, 
and otherwise $\trA$ may terminate.
%\commentt{Is ``high probability'' constant or higher?}
% and returns triangle free, and as before, this only adds negligible failure probability to $\mA'$. 
  This concludes the proof.
%%%
\iffalse
\paragraph{Adapting the proof to the case where $k$ is not divisible by 3.}
If $k$ is not divisible by $3$, then  to obtain $G^+$ from $G$, 
a similar process to the above is done with either $\lceil k/3\rceil, \lceil k/3\rceil, \lfloor k/3\rfloor$ or $\lceil k/3\rceil, \lfloor k/3\rfloor, \lfloor k/3\rfloor$, where all edges between the same two parts in the original graph $G$, are replaced by a path of the same length in $G^+$.
Then $G'$ is obtained from $G^+$ as before, and queries are answered as before, where the only difference is that the 

\paragraph{Degree queries $deg(u)$}
We can assume without loss of generality that no degree queries are performed, as $\mA'$ knows in advance the degree of every id. 
\commentt{This is different from what we had before. I think we should try to unify. I guess the first case should be made more similar to this one}

\paragraph{Neighbor queries $nbr(u, i)$}
If $u$ is an original vertex in $G'$, then $\mA$ queries $nbr(u,i)$ in $G$. It then performs a path labeling process so that each assigned id falls into the appropriate set of path vertices.  $\mA$ then returns the path.

If $u$ is a path vertex belonging to $V_{ij}$, then $\mA$ samples a uniform edge in $G$ between the relevant parts (explained below). Once an edge is sampled, $\mA$  performs a path labeling process and returns the path.
Sampling a uniform edge between $V_i, V_j$ is done by repeatedly choosing an index  $u\in V_i\cup V_j \setminus I_t$ and an index $r\in [d]$ and querying $nbr(u,r)$ in $G$ until the returned edge is between the relevant sets and was not previously witnessed.  If the edge sampling is preformed more than, say, $20\log n$, then $\mA$ aborts and returns that $G$ is triangle-free. 

\paragraph{pair queries} As before, we can assume w.l.o.g. that no pair queries are performed.

The rest of the algorithm is as before, except that if at any point the number of queries exceeds $C\cdot \sqrt{\trn}$ for some predetermined constant $C$, $\trA$ terminates and returns 
``triangle-free'' (or ``accept''). It remains to show that the query complexity and success probability do not change asymptotically due to the change in the answering process and potential early termination.

Let $q(t)$ denote the query complexity 
of the $t$th query. \commentt{DR: maybe clarify: the number of queries performed by $\mA$ on $G$ given the $t$th query of $\mA'$}
Observe that the success probability  of sampling an edge   between two specific parts that was not previously queried is at least $\frac{1}{2}-\frac{nd-|Q_n|}{nd}\geq 1/3$. \commentt{something here is a bit confusing with singular and plural}
Hence, 
the probability that in a single edge sampling process the algorithm performs more than $n^{0.1}$ queries is at most $(\frac{1}{3})^{20\log n}\leq \frac{1}{n^2}$.
By a union bound over all $Q_n$ queries, \commentt{of $\mA'$, right?}  with probability at least $1-1/n$, no single query is causing the algorithm to terminate early. We continue to show that with high probability the total  number of queries \commentt{performed by $\mA$} does not exceed $4Q_n$.
For any $t$,  $\Ex[q(t)]\leq 3$. 
Furthermore, for any $t$, $q(t)\leq 20\log n.$ 
Let $T=C\cdot \sqrt n$. By the additive Chernoff bound,
\[
 \Pr\left[\sum_{t=1}^{T} q(t)\geq 3Q_n+Q_n \right] < \exp\left( -\frac{Q_n^2}{Q_n\cdot 20\log n} \right)=o(1).
\]
Therefore, the probability that the algorithm terminates early is negligible, and so $\mA$ has the same asymptotic query complexity and success probability as does $\mA'$.
\fi
%
\end{proofof}
